\section{Experiments}

\subsection{Experimental Factors and Full Factorial Design}

Six factors are crossed fully (Table~\ref{tab:factors}), giving
\NumContexts{} contexts. Each factor exists to activate one of the policy
failure modes of Section~4.2: corridor demand and green ratio govern capacity
shortage; penetration governs how much of the demand a policy actually
commands, and hence the scope for herding; information lag governs staleness;
the incident regime supplies travel-time risk; and the alternative-capacity
factor determines whether there is uneven capacity to balance.

The demand levels span uncongested to oversaturated conditions relative to the
measured network capacity (Figure~\ref{fig:design}), and follow from that
measurement rather than from a chosen range.

\subsection{Common Random Numbers and Replications}

Each context is run with \NumSeeds{} paired seeds and all four policies, giving
\NumRunsMain{} runs in the main campaign.

The vehicle set, departure times, guided/unguided labelling, habitual corridor
choice, background traffic and the incident realisation depend on the seed and
the context alone, never on the policy. Two runs of a context under different
policies therefore face identical vehicles, and every comparison in this paper
is paired. Unguided drivers follow a habitual corridor drawn once from a logit
over free-flow path times fixed before any evaluation run; at penetration $p$, a
fraction $1-p$ of corridor demand is inert with respect to the policy under
test, which is what makes penetration a meaningful factor rather than a scaling
of demand.

\paragraph{Disruption as risk, not as a known event.} When the incident regime
is active, one lane closure is drawn --- location, onset and duration --- from a
declared distribution seeded by the run seed. Within a context and seed the
realisation is identical across all four policies, preserving the pairing;
across the seeds of a context it differs, so the context carries travel-time
risk that no policy can know in advance. An incident occurring at an identical
time and place in every replication would be a known event rather than a risk,
and could not test a reliability-aware policy at all.

\subsection{Main Campaign}

The main campaign is the full factorial of Section~5.1 at \NumSeeds{} seeds:
\NumRunsMain{} runs, from which every result in Section~\ref{sec:results}
derives unless the text says otherwise. Runs are validated against criteria
fixed in advance --- a completion rate of at least 0.98 and no teleported
vehicles --- and a context enters the analysis only if all four policies have
the full seed set.

\subsection{Distribution-Shift Experiments}

Eight held-out splits are evaluated separately and never pooled: demand above
and below the training range; an interior demand level as an interpolation
control; unseen information lag; unseen penetration; unseen green ratio;
training without disruption and testing under it; and a disruption located on a
corridor that carries none in training. Pooling them would average an
interpolation control together with a domain shift and hide the only thing the
experiment can say about them, which is that they differ.

The first seven splits are held out from the main campaign. The eighth requires
its own campaign, because no main-campaign context places a disruption on that
corridor; it comprises \NumRunsOODNorth{} additional runs.

\subsection{Boundary Refinement and Sensitivity Experiments}

Where the main grid shows the preferred policy changing between adjacent demand
levels, three intermediate demand levels were run in each of the \NumBoundN{}
widest-margin cells, a further \NumRunsBoundary{} runs. The purpose is to ask
whether the switching boundary can be located more finely than the
\NumBoundGrid~veh/h grid spacing.

Sensitivity to the two policy parameters fixed in advance --- the
load-balancing admissibility tolerance and the reliability dispersion weight ---
was measured on a declared subset of contexts, a further \NumRunsSens{} runs.
Sensitivity to the preference weighting is computed from the main campaign and
requires no additional simulation.

\subsection{Experimental Integrity and Protocol Deviations}
\label{sec:integrity}

\paragraph{What was fixed, and when.} The experiment specification --- factors
and levels, criteria, seed count, the resolution criterion of
Eq.~\ref{eq:resolved}, screening gates, baseline ladder, model class and
hyperparameters, abstention rules and shift splits --- was written and committed
to version control \emph{before any evaluation run was executed}. Policy
parameters were fixed at the same time and varied afterwards only as
sensitivity. No factor level was chosen because it produced a particular
outcome; the demand levels follow from measured capacity.

\paragraph{Execution reconciliation.} The main campaign comprises
\NumRunsMain{} valid runs
($\NumContexts \times \NumPolicies \times \NumSeeds$), and this is the number
underlying every result unless stated otherwise. The project executed
\NumRunsTotalAttempted{} simulations in total, of which
\NumRunsTotalSucceeded{} succeeded: \NumRunsValidate{} scenario-validation,
\NumRunsScreen{} screening, \NumRunsMain{} main, \NumRunsSens{} sensitivity,
\NumRunsOODNorth{} domain-shift and \NumRunsBoundary{} boundary-refinement runs,
plus \NumRunsFailedBypass{} runs that failed for the reason given below. These
counts are distinct and are not interchangeable.

\paragraph{Protocol deviation D-1.} The pre-specified complementarity screen did
not satisfy criterion G1 and therefore triggered the declared stopping rule. The
main campaign was nevertheless executed, under the documented protocol deviation
D-1. The reasons for continuation were recorded before campaign launch: a
selection problem demonstrably existed; G1 was written against the scalar
criterion and was blind to the three-criterion structure beside it; and the
negative finding rested on 16 of \NumContexts{} contexts sampled at extreme
factor levels only. A confirmatory criterion G1$'$, declared before the main
campaign and identical in definition to G1, was then evaluated on all
\NumContexts{} contexts. G1$'$ \NumGOneVerdict{}; the result is reported in
Section~\ref{sec:complementarity}.

We do not present G1$'$ as erasing the deviation, and we do not describe the
original screen as having passed. It did not. The methodological lesson is that
the screen sampled corner conditions and did not sample the interior transition
region, so it had power to detect near-ties but not coverage to find where the
minority policies win. A resolution-IV fraction at extreme levels is an
efficient design for estimating main effects and a poor one for discovering
where a winner map changes. That is a limitation of the screening design, not a
justification for having continued.

\paragraph{An infeasible shift split and its replacement.} One declared shift
split placed the disruption on the bypass corridor. The bypass is single-lane,
so the specified capacity-reducing lane closure does not constrain it --- it
disconnects it, leaving vehicles already routed there without a valid path. All
\NumRunsFailedBypass{} attempted runs failed for this reason. The split was
replaced by the nearest feasible construction, a disruption on the two-lane
north corridor, which preserves the intent --- a corridor carrying no disruption
in any training context --- while remaining a capacity reduction rather than a
severance. A guard now refuses any incident on a single-lane edge. The
replacement is reported as split O8 and is \emph{not} the originally declared
split.

\paragraph{Traceability.} Every numeric value reported in this paper is
generated from a single machine-readable results file by a script; a build-time
check fails if the manuscript references a quantity the analysis did not
produce. Automated quality-control checks cover leakage discipline,
training-fold-only baseline selection, freeze ordering, shift reporting and
citation completeness. All \NumRefs{} references were verified against publisher
records; DOIs that could not be confirmed directly were omitted rather than
inferred.
